\section{Tensors}

  % notation used only in this file
  \providecommand{\llbracket}{[\![}
  \providecommand{\rrbracket}{]\!]}

  Recall the definition of a \hyperref[lin-def:tensor]{tensor}. For our purposes, we will not need the full machinery of tensor product spaces, and it will be more useful to think of tensors concretely, as arrays of numbers.

  \begin{definition}[Tensor]
    \label{def:tensor}
    An \textbf{$N$-way tensor}, or \textbf{tensor of order $N$}, is a multidimensional array
    \begin{equation}
      \mathcal{X} \in \mathbb{R}^{I_1 \times I_2 \times \cdots \times I_N} \cong \mathbb{R}^{I_1} \otimes \mathbb{R}^{I_2} \otimes \cdots \otimes \mathbb{R}^{I_N}
    \end{equation}
    with entries $x_{i_1 i_2 \cdots i_N} \in \mathbb{R}$ for $1 \leq i_n \leq I_n$. Each index is called a \textbf{mode} (or way), and $I_n$ is the \textbf{dimension} of mode $n$.\footnote{In some fields, the order of the tensor is called its ``rank.'' We will not use this convention, since the word rank will mean something quite different in the next section.}
  \end{definition}

  Since we are working in $\mathbb{R}^n$ with the standard \hyperref[lin-def:inner-product]{inner product}, we will \textit{not} distinguish between $\mathbb{R}^n$ and its dual $(\mathbb{R}^n)^*$. That is, we won't care about covariant versus contravariant indices, and we will freely treat elements of $\mathbb{R}^n$ as \hyperref[lin-def:column-vector]{column vectors} or \hyperref[lin-def:row-vector]{row vectors}, whichever is more convenient. A vector is then a $1$-way tensor, and a matrix is a $2$-way tensor. Tensors of order $3$ or higher are called \textit{higher-order} tensors.

  \begin{example}[A $3$-Tensor]
    \label{ex:running}
    A $3$-way tensor $\mathcal{X} \in \mathbb{R}^{3 \times 4 \times 2}$ can be thought of as a box of numbers, and the easiest way to write it down is to list its two $3 \times 4$ \textbf{frontal slices} $X_k \coloneqq (x_{ijk})_{i,j}$.
    \begin{equation}
      X_1 = \begin{bmatrix} 1 & 4 & 7 & 10 \\ 2 & 5 & 8 & 11 \\ 3 & 6 & 9 & 12 \end{bmatrix}, \qquad
      X_2 = \begin{bmatrix} 13 & 16 & 19 & 22 \\ 14 & 17 & 20 & 23 \\ 15 & 18 & 21 & 24 \end{bmatrix}
    \end{equation}
    For instance, $x_{111} = 1$, $x_{231} = 8$, and $x_{342} = 24$. We will use this as our running example.
  \end{example}

\subsection{Rank-One Tensors and Tensor Rank}

  The basic building block of a tensor is the following.

  \begin{definition}[Outer Product]
    \label{def:outer-product}
    The \textbf{\hyperref[lin-def:outer-product]{outer product}} of vectors $a^{(1)} \in \mathbb{R}^{I_1}, \ldots, a^{(N)} \in \mathbb{R}^{I_N}$ is the $N$-way tensor $a^{(1)} \circ a^{(2)} \circ \cdots \circ a^{(N)} \in \mathbb{R}^{I_1 \times \cdots \times I_N}$ with entries
    \begin{equation}
      \big( a^{(1)} \circ a^{(2)} \circ \cdots \circ a^{(N)} \big)_{i_1 i_2 \cdots i_N} = a^{(1)}_{i_1} a^{(2)}_{i_2} \cdots a^{(N)}_{i_N}
    \end{equation}
    A tensor that can be written this way is called \textbf{rank-one}.
  \end{definition}

  For $N = 2$, $a \circ b = a b^T$ is the usual \hyperref[def:outer-product]{outer product} of vectors, which is exactly a rank-$1$ matrix. So we are just generalizing what we already know: every matrix of rank $D$ can be written as a sum of $D$ rank-one matrices, and the rank is the smallest such $D$. We take this to be the definition for tensors.

  \begin{definition}[Tensor Rank]
    \label{def:tensor-rank}
    The \textbf{rank} of a tensor $\mathcal{X}$, denoted $\mathrm{rank}(\mathcal{X})$, is the smallest $D$ such that $\mathcal{X}$ can be written as a sum of $D$ rank-one tensors,
    \begin{equation}
      \mathrm{rank}(\mathcal{X}) \coloneqq \min \bigg\{ D : \mathcal{X} = \sum_{d=1}^D a^{(1)}_d \circ a^{(2)}_d \circ \cdots \circ a^{(N)}_d \bigg\}
    \end{equation}
  \end{definition}

  This looks like a perfect analogue of matrix rank, but the two behave very differently. For a matrix, we can compute the rank with Gaussian elimination or the SVD in polynomial time. For tensors, \textit{there is no straightforward algorithm to determine the rank of a given tensor}; in fact, the problem is NP-hard. As a sign of how bad it is, there is a particular $9 \times 9 \times 9$ tensor whose rank is only known to lie between $18$ and $23$. In practice, the rank is determined numerically by trying to fit rank-$D$ models for various $D$. Some other differences are the following.

  \begin{enumerate}
    \item \textit{Rank depends on the field.} The rank of a real tensor can be different over $\mathbb{R}$ and $\mathbb{C}$, whereas for matrices it cannot.

    \begin{example}[Rank Over $\mathbb{R}$ vs. $\mathbb{C}$]
      \label{ex:rank-over-r-vs-c}
      Let $\mathcal{X} \in \mathbb{R}^{2 \times 2 \times 2}$ have frontal slices
      \begin{equation}
        X_1 = \begin{bmatrix} 1 & 0 \\ 0 & 1 \end{bmatrix}, \qquad X_2 = \begin{bmatrix} 0 & 1 \\ -1 & 0 \end{bmatrix}
      \end{equation}
      Then $\mathcal{X}$ has rank $3$ over $\mathbb{R}$ and rank $2$ over $\mathbb{C}$. Writing $\mathcal{X} = \sum_{d} a_d \circ b_d \circ c_d$ and stacking the vectors into the columns of factor matrices $A, B, C$, a rank decomposition over $\mathbb{R}$ is
      \begin{equation}
        A = \begin{bmatrix} 1 & 0 & 1 \\ 0 & 1 & -1 \end{bmatrix}, \quad
        B = \begin{bmatrix} 1 & 0 & 1 \\ 0 & 1 & 1 \end{bmatrix}, \quad
        C = \begin{bmatrix} 1 & 1 & 0 \\ -1 & 1 & 1 \end{bmatrix}
      \end{equation}
      whereas over $\mathbb{C}$ we can do it with only two components,
      \begin{equation}
        A = \frac{1}{\sqrt{2}} \begin{bmatrix} 1 & 1 \\ -i & i \end{bmatrix}, \quad
        B = \frac{1}{\sqrt{2}} \begin{bmatrix} 1 & 1 \\ i & -i \end{bmatrix}, \quad
        C = \begin{bmatrix} 1 & 1 \\ -i & i \end{bmatrix}
      \end{equation}
      The proof that the real rank is not $2$ is more involved (this example is due to Kruskal), and we omit it.
    \end{example}

    \item \textit{Maximum rank is not typical rank.} The \textbf{maximum rank} is the largest rank attainable, while a \textbf{typical rank} is any rank that occurs with positive probability (i.e. on a set of positive \hyperref[mt-def:lebesgue-measure]{Lebesgue measure}). For $I \times J$ matrices, both are equal to $\min\{I, J\}$. For tensors they can differ, and moreover over $\mathbb{R}$ there can be more than one typical rank (over $\mathbb{C}$ there is always exactly one).

    \begin{example}[Typical Ranks of $2 \times 2 \times 2$ Tensors]
      \label{ex:typical-ranks-of-2times2times2-tensors}
      If we draw every entry of a $2 \times 2 \times 2$ tensor iid from a standard normal, then Monte Carlo experiments show that about $79\%$ of them have rank $2$ and about $21\%$ have rank $3$. So both $2$ and $3$ are typical ranks over $\mathbb{R}$. Rank-one tensors are possible, but occur with probability zero.
    \end{example}

    \item \textit{Weak upper bounds.} For a general $\mathcal{X} \in \mathbb{R}^{I \times J \times K}$, the only known bound on the maximum rank is the weak one
    \begin{equation}
      \mathrm{rank}(\mathcal{X}) \leq \min\{IJ, IK, JK\}
    \end{equation}
    compared to the matrix case where it is simply $\min\{I, J\}$.
  \end{enumerate}

  Another important difference, which concerns the uniqueness of the decomposition, will be discussed after we introduce the CP decomposition.

\subsection{Structured Tensors}

  \begin{definition}[Cubical, Symmetric, Supersymmetric, Diagonal Tensors]
    \label{def:cubical-symmetric-supersymmetric-diagonal}
    Let $\mathcal{X} \in \mathbb{R}^{I_1 \times \cdots \times I_N}$.
    \begin{enumerate}
      \item $\mathcal{X}$ is \textbf{cubical} if every mode has the same size, i.e. $\mathcal{X} \in \mathbb{R}^{I \times I \times \cdots \times I}$.

      \item A cubical tensor is \textbf{supersymmetric}\footnote{Some authors prefer to just call this symmetric.} if its entries are invariant under any permutation of the indices. For example, $\mathcal{X} \in \mathbb{R}^{I \times I \times I}$ is supersymmetric if
      \begin{equation}
        x_{ijk} = x_{ikj} = x_{jik} = x_{jki} = x_{kij} = x_{kji}, \quad \forall i, j, k
      \end{equation}

      \item $\mathcal{X}$ can also be \textbf{symmetric} in just some of its modes, which requires those modes to have the same size. For example, $\mathcal{X} \in \mathbb{R}^{I \times I \times K}$ is symmetric in modes one and two if all of its frontal slices are symmetric,
      \begin{equation}
        X_k = X_k^T, \quad k = 1, \ldots, K
      \end{equation}

      \item $\mathcal{X}$ is \textbf{diagonal} if $x_{i_1 i_2 \cdots i_N} \neq 0$ only if $i_1 = i_2 = \cdots = i_N$.
    \end{enumerate}
  \end{definition}

  The entries $x_{ii \cdots i}$ of a cubical tensor make up its \textbf{superdiagonal}. An example is the cubical tensor with ones along the superdiagonal, which is diagonal and supersymmetric. For $N = 2$, these are the usual symmetric and diagonal matrices. Supersymmetric tensors are in bijection with homogeneous polynomials, much like \hyperref[lin-def:symmetric-matrices]{symmetric matrices} are in bijection with quadratic forms.

\subsection{Matricization}

  It is often useful to turn a tensor into a matrix so that we can use the linear algebra that we already know. This process is called \textbf{matricization}, \textbf{unfolding}, or \textbf{flattening}. In general, this can be any reordering of the entries of an $N$-way array into a matrix (e.g. a $2 \times 3 \times 4$ tensor can be arranged as a $6 \times 4$ or a $3 \times 8$ matrix), but we will only need a special case.

  A \textbf{mode-$n$ fiber} of $\mathcal{X}$ is a vector obtained by fixing every index except the $n$th, which is the higher-order analogue of the rows and columns of a matrix. For a $3$-way tensor, the mode-$1$, mode-$2$, mode-$3$ fibers are the columns $x_{:jk}$, the rows $x_{i:k}$, and the tubes $x_{ij:}$, respectively.

  \begin{definition}[Mode-$n$ Matricization]
    \label{def:mode-n-matricization}
    The \textbf{mode-$n$ matricization} of $\mathcal{X} \in \mathbb{R}^{I_1 \times \cdots \times I_N}$ is the matrix $X_{(n)} \in \mathbb{R}^{I_n \times (I_1 \cdots I_{n-1} I_{n+1} \cdots I_N)}$ whose columns are the mode-$n$ fibers of $\mathcal{X}$. Precisely, the entry $x_{i_1 i_2 \cdots i_N}$ is mapped to the matrix entry $(i_n, j)$ where
    \begin{equation}
      j = 1 + \sum_{\substack{k=1 \\ k \neq n}}^N (i_k - 1) J_k, \qquad J_k = \prod_{\substack{m=1 \\ m \neq n}}^{k-1} I_m
    \end{equation}
  \end{definition}

  The formal notation is clunky, so let's just do an example.

  \begin{example}[Unfoldings of a $3 \times 4 \times 2$ Tensor]
    \label{ex:unfoldings-of-a-3times4times2-tensor}
    Let $\mathcal{X} \in \mathbb{R}^{3 \times 4 \times 2}$ be the tensor from Example \ref{ex:running}. Then the three mode-$n$ unfoldings are
    \begin{align}
      X_{(1)} & = \begin{bmatrix} 1 & 4 & 7 & 10 & 13 & 16 & 19 & 22 \\ 2 & 5 & 8 & 11 & 14 & 17 & 20 & 23 \\ 3 & 6 & 9 & 12 & 15 & 18 & 21 & 24 \end{bmatrix} \in \mathbb{R}^{3 \times 8} \\
      X_{(2)} & = \begin{bmatrix} 1 & 2 & 3 & 13 & 14 & 15 \\ 4 & 5 & 6 & 16 & 17 & 18 \\ 7 & 8 & 9 & 19 & 20 & 21 \\ 10 & 11 & 12 & 22 & 23 & 24 \end{bmatrix} \in \mathbb{R}^{4 \times 6} \\
      X_{(3)} & = \begin{bmatrix} 1 & 2 & 3 & 4 & \cdots & 10 & 11 & 12 \\ 13 & 14 & 15 & 16 & \cdots & 22 & 23 & 24 \end{bmatrix} \in \mathbb{R}^{2 \times 12}
    \end{align}
    So $X_{(1)} = [X_1 \;\; X_2]$ is just the frontal slices placed side by side.
  \end{example}

  Different papers order the columns differently, but the specific permutation of the columns is not important as long as we are consistent. Similarly, we can \textbf{vectorize} a tensor into $\mathrm{vec}(\mathcal{X}) \in \mathbb{R}^{I_1 \cdots I_N}$, and in the example above $\mathrm{vec}(\mathcal{X}) = (1, 2, \ldots, 24)^T$.

\subsection{The $n$-Mode Product}

  Now we want to multiply tensors. There are many ways of doing this, but we will only consider multiplying a tensor by a matrix (or vector) along a single mode.

  \begin{definition}[$n$-Mode Product]
    \label{def:n-mode-product}
    Let $\mathcal{X} \in \mathbb{R}^{I_1 \times \cdots \times I_N}$.
    \begin{enumerate}
      \item The \textbf{$n$-mode (matrix) product} of $\mathcal{X}$ with a matrix $U \in \mathbb{R}^{J \times I_n}$ is the tensor $\mathcal{X} \times_n U \in \mathbb{R}^{I_1 \times \cdots \times I_{n-1} \times J \times I_{n+1} \times \cdots \times I_N}$ with entries
      \begin{equation}
        (\mathcal{X} \times_n U)_{i_1 \cdots i_{n-1} \, j \, i_{n+1} \cdots i_N} = \sum_{i_n = 1}^{I_n} x_{i_1 i_2 \cdots i_N} \, u_{j i_n}
      \end{equation}

      \item The \textbf{$n$-mode (vector) product} of $\mathcal{X}$ with a vector $v \in \mathbb{R}^{I_n}$ is the tensor $\mathcal{X} \,\bar{\times}_n\, v \in \mathbb{R}^{I_1 \times \cdots \times I_{n-1} \times I_{n+1} \times \cdots \times I_N}$ of order $N - 1$ with entries
      \begin{equation}
        (\mathcal{X} \,\bar{\times}_n\, v)_{i_1 \cdots i_{n-1} \, i_{n+1} \cdots i_N} = \sum_{i_n = 1}^{I_n} x_{i_1 i_2 \cdots i_N} \, v_{i_n}
      \end{equation}
    \end{enumerate}
  \end{definition}

  In the matrix case, each mode-$n$ fiber of $\mathcal{X}$ gets multiplied by $U$, and in the vector case, we take the \hyperref[lin-def:inner-product]{inner product} of each mode-$n$ fiber with $v$. In terms of unfoldings, this means that
  \begin{equation}
    \mathcal{Y} = \mathcal{X} \times_n U \iff Y_{(n)} = U X_{(n)}
  \end{equation}
  When $\mathcal{X}$ defines a multilinear operator, the $n$-mode product is just a change of basis in mode $n$.

  Now, is this something we've already seen? Yes. Since we are not distinguishing between $\mathbb{R}^n$ and its dual, summing over a pair of shared indices is exactly a \textbf{tensor contraction}. The $n$-mode product $\mathcal{X} \times_n U$ is the contraction of $\mathcal{X}$ and $U$ along the index $i_n$ of $\mathcal{X}$ and the column index of $U$, followed by moving the leftover row index $j$ of $U$ into position $n$. Similarly, $\mathcal{X} \,\bar{\times}_n\, v$ is the contraction of $\mathcal{X}$ with $v$ along mode $n$, which is why the order drops by one.

  \begin{example}[Mode Products of the Running Tensor]
    \label{ex:mode-products-of-the-running-tensor}
    Let $\mathcal{X}$ be as in Example \ref{ex:running}.
    \begin{enumerate}
      \item Let $U = \begin{bmatrix} 1 & 3 & 5 \\ 2 & 4 & 6 \end{bmatrix}$. Then $\mathcal{Y} = \mathcal{X} \times_1 U \in \mathbb{R}^{2 \times 4 \times 2}$ has frontal slices $Y_k = U X_k$,
      \begin{equation}
        Y_1 = \begin{bmatrix} 22 & 49 & 76 & 103 \\ 28 & 64 & 100 & 136 \end{bmatrix}, \qquad
        Y_2 = \begin{bmatrix} 130 & 157 & 184 & 211 \\ 172 & 208 & 244 & 280 \end{bmatrix}
      \end{equation}

      \item Let $v = (1, 2, 3, 4)^T$. Then
      \begin{equation}
        \mathcal{X} \,\bar{\times}_2\, v = \begin{bmatrix} 70 & 190 \\ 80 & 200 \\ 90 & 210 \end{bmatrix} \in \mathbb{R}^{3 \times 2}
      \end{equation}
      For instance, the top-left entry is $\langle x_{1:1}, v \rangle = 1 + 8 + 21 + 40 = 70$.
    \end{enumerate}
  \end{example}

\subsection{Matrix Products}

  Several matrix products show up when relating tensors to their unfoldings, so let's define them.

  \begin{definition}[Kronecker, Khatri-Rao, Hadamard Products]
    \label{def:kronecker-khatri-rao-hadamard-products}
    \begin{enumerate}
      \item The \textbf{\hyperref[lin-def:kronecker-product]{Kronecker product}} of $A \in \mathbb{R}^{I \times J}$ and $B \in \mathbb{R}^{K \times L}$ is the $(IK) \times (JL)$ matrix
      \begin{equation}
        A \otimes B = \begin{bmatrix} a_{11} B & a_{12} B & \cdots & a_{1J} B \\ a_{21} B & a_{22} B & \cdots & a_{2J} B \\ \vdots & \vdots & \ddots & \vdots \\ a_{I1} B & a_{I2} B & \cdots & a_{IJ} B \end{bmatrix} = \begin{bmatrix} a_1 \otimes b_1 & a_1 \otimes b_2 & \cdots & a_J \otimes b_L \end{bmatrix}
      \end{equation}

      \item The \textbf{Khatri-Rao product} is the ``matching columnwise'' Kronecker product. For $A \in \mathbb{R}^{I \times K}$ and $B \in \mathbb{R}^{J \times K}$, it is the $(IJ) \times K$ matrix
      \begin{equation}
        A \odot B = \begin{bmatrix} a_1 \otimes b_1 & a_2 \otimes b_2 & \cdots & a_K \otimes b_K \end{bmatrix}
      \end{equation}

      \item The \textbf{Hadamard product} is the elementwise product. For $A, B \in \mathbb{R}^{I \times J}$, it is the $I \times J$ matrix $(A * B)_{ij} = a_{ij} b_{ij}$.
    \end{enumerate}
  \end{definition}

  If $a, b$ are vectors, then the Khatri-Rao and \hyperref[lin-def:kronecker-product]{Kronecker products} coincide, $a \odot b = a \otimes b$. These products, along with the $n$-mode product, satisfy the following properties, which we will use over and over.

  \begin{lemma}[Properties of Tensor and Matrix Products]
    \label{thm:products}
    Let $A^\dagger$ denote the Moore-Penrose pseudoinverse of $A$.
    \begin{enumerate}
      \item \textit{(Kronecker)} For matrices of compatible sizes,
      \begin{align}
        (A \otimes B)(C \otimes D) & = AC \otimes BD \\
        (A \otimes B)^\dagger & = A^\dagger \otimes B^\dagger
      \end{align}

      \item \textit{(Khatri-Rao)} The Khatri-Rao product is associative,
      \begin{equation}
        A \odot B \odot C = (A \odot B) \odot C = A \odot (B \odot C)
      \end{equation}
      and,
      \begin{align}
        (A \odot B)^T (A \odot B) & = A^T A * B^T B \\
        (A \odot B)^\dagger & = \big( (A^T A) * (B^T B) \big)^\dagger (A \odot B)^T
      \end{align}

      \item \textit{($n$-mode, distinct modes)} The order of multiplication in distinct modes is irrelevant,
      \begin{equation}
        \mathcal{X} \times_m A \times_n B = \mathcal{X} \times_n B \times_m A \qquad (m \neq n)
      \end{equation}

      \item \textit{($n$-mode, same mode)}
      \begin{equation}
        \mathcal{X} \times_n A \times_n B = \mathcal{X} \times_n (BA)
      \end{equation}

      \item \textit{($n$-mode, vectors)} For vectors, precedence matters since the modes get renumbered after each product. For $m < n$,
      \begin{equation}
        \mathcal{X} \,\bar{\times}_m\, a \,\bar{\times}_n\, b = (\mathcal{X} \,\bar{\times}_m\, a) \,\bar{\times}_{n-1}\, b = (\mathcal{X} \,\bar{\times}_n\, b) \,\bar{\times}_m\, a
      \end{equation}

      \item \textit{(Unfolding)} Let $\mathcal{X} \in \mathbb{R}^{I_1 \times \cdots \times I_N}$ and $A^{(n)} \in \mathbb{R}^{J_n \times I_n}$ for all $n$. Then for any $n$,
      \begin{equation}
        \mathcal{Y} = \mathcal{X} \times_1 A^{(1)} \times_2 A^{(2)} \cdots \times_N A^{(N)} \iff Y_{(n)} = A^{(n)} X_{(n)} \big( A^{(N)} \otimes \cdots \otimes A^{(n+1)} \otimes A^{(n-1)} \otimes \cdots \otimes A^{(1)} \big)^T
      \end{equation}
    \end{enumerate}
  \end{lemma}
  \begin{proof}
    The Kronecker and Khatri-Rao identities are direct computations. For (3) and (4), note that by the unfolding relation $Y_{(n)} = U X_{(n)}$, multiplying by $A$ and then $B$ in mode $n$ gives $B (A X_{(n)}) = (BA) X_{(n)}$, which is (4). Distinct modes act on different indices, so the sums can be exchanged, which is (3). The proof of (6) can be found in Kolda's paper on multilinear operators; TBD.
  \end{proof}

\subsection{CP Decomposition}

  We can now state the first of the two main tensor decompositions. The idea is to write a tensor as a sum of rank-one tensors, and in this sense it generalizes the SVD, where $X = \sum_d \sigma_d u_d v_d^T$. It goes by many names: the polyadic form of a tensor (Hitchcock, 1927), PARAFAC for parallel factors (Harshman, 1970), CANDECOMP for canonical decomposition (Carroll and Chang, 1970), and now just \textbf{CP} for CANDECOMP/PARAFAC.

  \begin{definition}[CP Decomposition]
    \label{def:cp-decomposition}
    A \textbf{CP decomposition} of $\mathcal{X} \in \mathbb{R}^{I \times J \times K}$ with $D$ components is an approximation
    \begin{equation}
      \mathcal{X} \approx \sum_{d=1}^D a_d \circ b_d \circ c_d \eqqcolon \llbracket A, B, C \rrbracket
    \end{equation}
    where $a_d \in \mathbb{R}^I$, $b_d \in \mathbb{R}^J$, $c_d \in \mathbb{R}^K$. Elementwise, this is $x_{ijk} \approx \sum_{d=1}^D a_{id} b_{jd} c_{kd}$. The matrices $A = [a_1 \; \cdots \; a_D] \in \mathbb{R}^{I \times D}$, and likewise $B \in \mathbb{R}^{J \times D}$, $C \in \mathbb{R}^{K \times D}$, are called the \textbf{factor matrices}. More generally, for $\mathcal{X} \in \mathbb{R}^{I_1 \times \cdots \times I_N}$,
    \begin{equation}
      \mathcal{X} \approx \llbracket \lambda; A^{(1)}, \ldots, A^{(N)} \rrbracket \coloneqq \sum_{d=1}^D \lambda_d \, a^{(1)}_d \circ a^{(2)}_d \circ \cdots \circ a^{(N)}_d
    \end{equation}
    where $A^{(n)} \in \mathbb{R}^{I_n \times D}$ has columns normalized to unit length and the weights are absorbed into $\lambda \in \mathbb{R}^D$.
  \end{definition}

  If the approximation is an equality and $D = \mathrm{rank}(\mathcal{X})$, then this is called a \textbf{rank decomposition}. Using the matricization and Khatri-Rao products above, the \hyperref[def:cp-decomposition]{CP decomposition} can be written in matrix form, one for each mode.
  \begin{align}
    X_{(1)} & \approx A (C \odot B)^T \\
    X_{(2)} & \approx B (C \odot A)^T \\
    X_{(3)} & \approx C (B \odot A)^T
  \end{align}
  and for general $N$, with $\Lambda = \mathrm{diag}(\lambda)$,
  \begin{equation}
    X_{(n)} \approx A^{(n)} \Lambda \big( A^{(N)} \odot \cdots \odot A^{(n+1)} \odot A^{(n-1)} \odot \cdots \odot A^{(1)} \big)^T
  \end{equation}

  \subsubsection{Uniqueness}

    Here is one more major difference between matrices and tensors, and this time it's in the tensors' favor. Matrix decompositions are not unique. If $X = A B^T$ has rank $D$, then for any orthogonal $W \in \mathbb{R}^{D \times D}$ we also have $X = (AW)(BW)^T$. The SVD is unique (when the singular values are distinct) only because we add orthogonality constraints and order the singular values. The CP decomposition, on the other hand, is unique under much weaker conditions, up to two trivial indeterminacies.
    \begin{enumerate}
      \item \textit{Permutation.} The rank-one components can be reordered, $\llbracket A, B, C \rrbracket = \llbracket A \Pi, B \Pi, C \Pi \rrbracket$ for any \hyperref[lin-def:permutation-matrix]{permutation matrix} $\Pi$.
      \item \textit{Scaling.} We can rescale each component, $\sum_d (\alpha_d a_d) \circ (\beta_d b_d) \circ (\gamma_d c_d)$, as long as $\alpha_d \beta_d \gamma_d = 1$ for all $d$.
    \end{enumerate}

    \begin{theorem}[Kruskal]
      \label{thm:kruskal}
      Let $\mathcal{X} = \llbracket A, B, C \rrbracket$ with $D$ components. Let the \textbf{$k$-rank} $k_A$ of a matrix $A$ be the largest $k$ such that any $k$ columns of $A$ are linearly independent. If
      \begin{equation}
        k_A + k_B + k_C \geq 2D + 2
      \end{equation}
      then the decomposition is unique up to permutation and scaling.
    \end{theorem}

\subsection{Tucker Decomposition}

  The second decomposition is a higher-order version of PCA. Instead of a sum of rank-one tensors, it writes the tensor as a small \textit{core} tensor that is transformed by a matrix along each mode. It is also called three-mode factor analysis, three-mode PCA, $N$-mode PCA, or higher-order SVD (HOSVD).

  \begin{definition}[Tucker Decomposition]
    \label{def:tucker-decomposition}
    A \textbf{Tucker decomposition} of $\mathcal{X} \in \mathbb{R}^{I \times J \times K}$ is an approximation
    \begin{equation}
      \mathcal{X} \approx \mathcal{G} \times_1 A \times_2 B \times_3 C = \sum_{p=1}^{R_1} \sum_{q=1}^{R_2} \sum_{r=1}^{R_3} g_{pqr} \, a_p \circ b_q \circ c_r \eqqcolon \llbracket \mathcal{G}; A, B, C \rrbracket
    \end{equation}
    where $\mathcal{G} \in \mathbb{R}^{R_1 \times R_2 \times R_3}$ is the \textbf{core tensor}, and $A \in \mathbb{R}^{I \times R_1}$, $B \in \mathbb{R}^{J \times R_2}$, $C \in \mathbb{R}^{K \times R_3}$ are the \textbf{factor matrices}, which are usually taken to have orthonormal columns. For $\mathcal{X} \in \mathbb{R}^{I_1 \times \cdots \times I_N}$,
    \begin{equation}
      \mathcal{X} \approx \mathcal{G} \times_1 A^{(1)} \times_2 A^{(2)} \cdots \times_N A^{(N)} = \llbracket \mathcal{G}; A^{(1)}, \ldots, A^{(N)} \rrbracket
    \end{equation}
    with $\mathcal{G} \in \mathbb{R}^{R_1 \times \cdots \times R_N}$ and $A^{(n)} \in \mathbb{R}^{I_n \times R_n}$.
  \end{definition}

  The columns of the factor matrices can be thought of as principal components in each mode, and the entries of the core $\mathcal{G}$ describe the level of interaction between the components of different modes. If $R_n < I_n$, then $\mathcal{G}$ is a compressed version of $\mathcal{X}$, and the storage needed for $(\mathcal{G}, A^{(1)}, \ldots, A^{(N)})$ can be much smaller than for $\mathcal{X}$. By the Unfolding property in Lemma \ref{thm:products}, the matricized forms are
  \begin{equation}
    X_{(n)} \approx A^{(n)} G_{(n)} \big( A^{(N)} \otimes \cdots \otimes A^{(n+1)} \otimes A^{(n-1)} \otimes \cdots \otimes A^{(1)} \big)^T
  \end{equation}

  \begin{lemma}[CP is a Special Case of Tucker]
    \label{thm:cp-is-a-special-case-of-tucker}
    The \hyperref[def:cp-decomposition]{CP decomposition} $\llbracket A, B, C \rrbracket$ with $D$ components is the \hyperref[def:tucker-decomposition]{Tucker decomposition} with $R_1 = R_2 = R_3 = D$ and a diagonal core $\mathcal{G}$ with ones on the superdiagonal.
  \end{lemma}
  \begin{proof}
    With $g_{pqr} = 1$ if $p = q = r$ and $0$ otherwise, the triple sum collapses to $\sum_{d} a_d \circ b_d \circ c_d$.
  \end{proof}

  Just as the rank of a matrix is the dimension of its \hyperref[lin-def:column-space]{column space}, we can define the following for Tucker.

  \begin{definition}[$n$-Rank]
    \label{def:n-rank}
    The \textbf{$n$-rank} of $\mathcal{X}$, denoted $\mathrm{rank}_n(\mathcal{X})$, is the column rank of $X_{(n)}$, i.e. the dimension of the span of the mode-$n$ fibers. If $R_n = \mathrm{rank}_n (\mathcal{X})$, then $\mathcal{X}$ is a \textbf{rank-$(R_1, \ldots, R_N)$} tensor.
  \end{definition}

  Trivially $R_n \leq I_n$. Note that this is \textit{not} the same as the \hyperref[def:tensor-rank]{tensor rank} defined earlier. It is always easy to find an exact \hyperref[def:tucker-decomposition]{Tucker decomposition} of $\mathcal{X}$ with $R_n = \mathrm{rank}_n(\mathcal{X})$. If we pick some $R_n < \mathrm{rank}_n(\mathcal{X})$, then the decomposition is necessarily inexact (a \textit{truncated} Tucker decomposition) and harder to compute.

  Finally, we can restrict some of the factor matrices to be the identity. The \textbf{Tucker2} decomposition sets one factor matrix to the identity, e.g. $\mathcal{X} \approx \mathcal{G} \times_1 A \times_2 B = \llbracket \mathcal{G}; A, B, I \rrbracket$, and \textbf{Tucker1} sets two of them, e.g. $\mathcal{X} \approx \mathcal{G} \times_1 A$, which has $X_{(1)} \approx A G_{(1)}$ and is therefore equivalent to ordinary PCA on the mode-$1$ unfolding.
